\section{Towards A Better One}
\subsection{History of Load Reversals}
In conventional plasticity models, local reversals cannot be easily detected.
A common criterion for detecting reversal is the double contraction $\beeta:\dot{\bs}$.
For example, in the multi-surface theory \citep{Semenov2022}, the exact criterion is used to identify the initiation of a load reversal, which further controls the activation of a new surface.

The normal yield ratio $z$ effectively represents the current `loading level' at any given moment.
When it increases, the subloading surface expands, indicating a loading process.
When it decreases, the subloading surface contracts, indicating an unloading process.
Thus, with the introduction of $z$, local reversals can be detected with ease by checking whether $z$ ever decreases during the given (sub)step.

Let the normal yield ratio at each load reversal be denoted as $z_r^j$ with $j=1,2,3,\cdots$ standing for the $j$-th reversal.
The initial value is $z_r^0=0$.
We use $z_r$ to denote the collection of all $z_r^j$, thus,
\begin{gather}
    z_r=\{z_r^j\}.
\end{gather}